% ============================================================
% Chapter 4: Proposed Method: Specification and Design
%   Section: System Architecture
% Included from chapters/04-proposed-method/proposed-method.tex
% Sources: docs/arquitectura-tecnica.md section 2.1 (topology
% diagram), docs/decisions/inference.md, docker/docker-compose.yml.
% ============================================================

\section{System Architecture}
\label{sec:architecture}

The system is a set of services run together by Docker Compose
(Figure~\ref{fig:architecture}).

\begin{itemize}
    \item \textbf{Backend} (Python, FastAPI): the API and all the
    business logic --- discovery, selection, extraction, enrichment,
    admission, fact-checking, storage and the evaluation harness. It
    holds the vector store (Qdrant, embedded in the process) and the
    analysis cache, and journals every job event to disk.
    \item \textbf{Inference service} (Python, FastAPI): a model server
    for the three transformer models --- GLiNER for entities, an
    XLM-RoBERTa sentiment classifier and the BGE-M3 embedding model.
    It knows nothing about topics, thresholds or scoring; the backend
    reaches it over HTTP.
    \item \textbf{SearXNG}: a self-hosted metasearch engine that sends
    each query to a closed list of web and academic search engines and
    merges their answers.
    \item \textbf{Ollama}: the language model, behind an
    OpenAI-compatible endpoint, so a hosted provider can replace it by
    configuration alone.
    \item \textbf{Neo4j}: the graph database every analysis is written
    to.
    \item \textbf{Frontend} (Next.js, TypeScript): the newsroom and lab
    pages and the public reader view. It calls the backend from the
    server side only, so the API key never reaches a browser.
\end{itemize}

\begin{figure}[ht]
\centering
\begin{adjustbox}{max width=\linewidth}
\begin{tikzpicture}[
    node distance=9mm and 12mm,
    svc/.style={draw, rounded corners=2pt, align=center, font=\footnotesize,
                minimum height=9mm, text width=3.1cm, fill=ganttplan!15},
    ext/.style={svc, fill=black!6},
    db/.style={svc, fill=ganttdone!25},
    arr/.style={-{Stealth[length=2mm]}, thick},
    lbl/.style={font=\scriptsize\itshape, fill=white, inner sep=1pt},
]
\node[svc] (fe) {Frontend\\Next.js};
\node[svc, right=of fe, text width=3.6cm, minimum height=16mm] (be)
    {Backend (FastAPI)\\discovery, selection, extraction,
     enrichment, admission, fact-checking, storage};
\node[svc, right=of be] (inf) {Inference service\\GLiNER, XLM-R, BGE-M3};
\node[ext, below=of be, xshift=-2.6cm] (sx) {SearXNG\\web + academic engines};
\node[ext, below=of be, xshift=1.1cm] (llm) {Ollama\\\texttt{llama3.2:3b}};
\node[db, below=of inf] (neo) {Neo4j\\graph};
\node[db, above=of be, xshift=-2.2cm] (lake) {Data lake\\raw / processed / exploitation};
\node[db, above=of be, xshift=2.2cm] (qd) {Qdrant\\(embedded)};
\node[ext, left=of sx] (web) {Open web and\\configured sources};

\draw[arr] (fe) -- node[lbl, above]{HTTP} (be);
\draw[arr] (be) -- node[lbl, above]{HTTP} (inf);
\draw[arr] (be) -- (sx);
\draw[arr] (be) -- (llm);
\draw[arr] (be) -- node[lbl, sloped, above]{fail-soft} (neo);
\draw[arr] (be) -- (lake);
\draw[arr] (be) -- (qd);
\draw[arr] (sx) -- (web);
\draw[arr] (be.west) to[bend right=20] node[lbl, near end, left]{fetch, guarded} (web.north);
\end{tikzpicture}
\end{adjustbox}
\caption{The services and how they talk. Every arrow from the backend
to an external service passes through that service's concurrency
ceiling (Section~\ref{sec:design-concurrency}).}
\label{fig:architecture}
\end{figure}

\paragraph{Why the models live in their own service.} The three
transformer models used to load inside the backend, beside the API. The
split was made for three measured reasons: the machine-learning wheels
dominated and broke the backend's image build; models loaded lazily on
the first request, so the first analysis after a restart silently paid
a 10--15~s load; and the models measure 3.58~GiB resident, while the
backend container was capped at 1~GiB, which made an out-of-memory kill
on the first real analysis likely. After the split the backend runs in
about 230~MiB and the inference service in 3.58~GiB; the inference
service reports itself healthy only once all three models are loaded,
and the backend waits for it. A test of the repository fails if the
backend ever imports a model library again.

\paragraph{The language model is a service too.} The verifier, the
selection step and the text corrector all reach the model through one
client that speaks the OpenAI wire format. Switching from the local
model to a hosted provider is a change of two settings, which is what
makes the model benchmarks of Phase~5 one command per model.
